% !TeX program = xelatex
\documentclass[11pt]{article}
\newcommand{\notesdir}{../../lecture-notes}
\providecommand{\notesdir}{.}

% -----------------------------------------------------------------------------
% Shared preamble for the CO5059 lecture notes and the quiz-review documents.
% Documents outside this folder define \notesdir (path to this folder) first.
% -----------------------------------------------------------------------------

% Page layout and typography
\usepackage[a4paper,portrait,margin=1in]{geometry}
\usepackage{fontspec}
\newfontfamily\arialfont{Arial}
\usepackage{microtype}
\usepackage{parskip}

% Mathematics and diagrams
\usepackage{amsmath,amssymb,amsthm,mathtools}
\usepackage{tikz-cd}
\usetikzlibrary{arrows.meta,positioning,shapes.geometric,shapes.arrows,calc,fit,backgrounds}

% Reusable styles for flowcharts and concept diagrams.
\tikzset{
  concept/.style={
    draw=brandprimary,
    fill=brandsecondary!8,
    rounded corners,
    minimum width=3.2cm,
    minimum height=0.9cm,
    align=center,
    font=\small,
  },
  diagram arrow/.style={-{Latex[length=2.5mm]},thick,brandprimary},
}

% Tables, figures, code, and lists
\usepackage{graphicx}
\graphicspath{{\notesdir/}}
\usepackage{enumitem}
\usepackage{booktabs}
\usepackage{array}
\usepackage{listings}

% Color and links
\usepackage[dvipsnames]{xcolor}
\definecolor{brandprimary}{HTML}{030391}
\definecolor{brandsecondary}{HTML}{1488D8}
\usepackage[
  colorlinks=true,
  linkcolor=brandprimary,
  citecolor=brandsecondary,
  urlcolor=brandsecondary
]{hyperref}

% Headers and footers
\usepackage{fancyhdr}
\setlength{\headheight}{36pt}
\pagestyle{fancy}
\fancyhf{}
\fancyhead[L]{%
  \parbox[c][0.42in][c]{0.42\textwidth}{%
    \raggedright\small\textbf{\coursecode}\\\coursename}}
\fancyhead[R]{%
  \parbox[c][0.42in][c]{0.42\textwidth}{%
    \raggedleft\includegraphics[height=0.42in]{assets/hcmut_logo.png}}}
\fancyfoot[L]{\small Semester \semester}
\fancyfoot[C]{\small\nouppercase{\leftmark}}
\fancyfoot[R]{\small Page \thepage}
\renewcommand{\sectionmark}[1]{\markboth{#1}{}}

% -----------------------------------------------------------------------------
% Course information -- change these values once per course
% -----------------------------------------------------------------------------
\newcommand{\universityname}{HCMC University of Technology}
\newcommand{\faculty}{Faculty of Computer Science and Engineering}
\newcommand{\coursecode}{CO5059}
\newcommand{\coursename}{Advanced System Architectures}
\newcommand{\semester}{261}
\newcommand{\instructor}{Tran Ngoc Thinh}
\newcommand{\instructoremail}{}% leave empty to hide
\newcommand{\notetaker}{Phan Nhat Minh}
\newcommand{\notetakeremail}{minh.phannhat95@hcmut.edu.vn}
\newcommand{\coursedescription}{Quantitative design and evaluation of computer
systems: performance and cost, instruction set architectures, pipelining and
its hazards, memory hierarchy, multiprocessors, superscalar/VLIW processors,
and AI \& security chips.}
\newcommand{\emaillink}[1]{%
  \href{mailto:#1}{\textcolor{brandprimary}{\texttt{#1}}}}

% Colored definitions, results, examples, proofs, and notes.
\input{\notesdir/colored-boxes}

% -----------------------------------------------------------------------------
% Course-specific notation
% -----------------------------------------------------------------------------
\newcommand{\CPI}{\mathrm{CPI}}

\lstdefinestyle{mips}{
  basicstyle=\ttfamily\small,
  keywordstyle=\bfseries\color{brandprimary},
  morekeywords={lw,sw,add,sub,and,or,slt,beq,bne,beqz,xor,nop,addi,addiu,addu,
    subu,li,la,jal,jr,jalr,sll,srl,sra,sllv,srlv,srav,andi,ori,xori,nor,slti,
    sltiu,sltu,lui,mult,multu,div,divu,mfhi,mflo,ll,sc,syscall,move,lh,lhu,lb,
    lbu,sh,sb,blez,bgtz,bltz,bgez,j,mul},
  showstringspaces=false,
  columns=fullflexible,
  keepspaces=true,
}
\lstset{style=mips}

% Creates a consistently formatted lecture heading.
% Usage: \lecture{number}{title}{date}
\newcommand{\lecture}[3]{%
  \section{Lecture #1: #2}
  \label{lec:#1}
  \begin{center}
    \small\textit{#3}
  \end{center}
}


% Cover page + contents page.  Usage: \makecover{Big title}{Description}
\newcommand{\makecover}[2]{%
\pagenumbering{roman}
\begin{titlepage}
  \thispagestyle{empty}

  \begin{center}
    \includegraphics[width=0.82\textwidth]{assets/hcmut_banner_en.png}

    \vspace{0.35cm}
    {\arialfont\large\bfseries\color{brandprimary}\faculty\par}
  \end{center}

  \vfill
  \begin{center}
    {\Huge\bfseries #1\par}
    \vspace{0.55cm}
    {\LARGE\bfseries\color{brandprimary}\coursename\par}
    \vspace{0.2cm}
    {\arialfont\large\bfseries\color{brandsecondary}\coursecode\par}

    \vspace{0.65cm}
    {\color{brandsecondary}\rule{0.38\textwidth}{1.2pt}}
    \vspace{0.65cm}

    {\large\textbf{Instructor:} \instructor
      \ifx\instructoremail\empty\else\ {\small(\emaillink{\instructoremail})}\fi\par}
    \vspace{0.3cm}
    {\large\textbf{Notes by:} \notetaker\
      {\small(\emaillink{\notetakeremail})}\par}
    \vspace{0.3cm}
    {\large\textbf{Semester:} \semester\par}

    \vspace{0.8cm}
    \begin{minipage}{0.72\textwidth}
      \centering\itshape\color{black!70}
      #2
    \end{minipage}
  \end{center}

  \vfill
  \begin{center}
    \small\color{black!60}Last updated: \today
  \end{center}
\end{titlepage}

% -----------------------------------------------------------------------------
% Table of contents -- kept on its own page
% -----------------------------------------------------------------------------
\clearpage
\setcounter{page}{2}
\thispagestyle{empty}
\hypersetup{linkcolor=black}
\setcounter{tocdepth}{2}
\tableofcontents
\hypersetup{linkcolor=brandprimary}
\clearpage
\pagenumbering{arabic}
}


\newcommand{\ex}[3]{\subsection*{Exercise #1 \normalfont\small(book p.~#2, PDF p.~#3)}%
  \addcontentsline{toc}{subsection}{Exercise #1}}

\begin{document}
\sloppy
\makecover{Textbook Exercises: Instruction Sets}{Hennessy \& Patterson,
\emph{Computer Architecture: A Quantitative Approach}, 5th ed., Appendix~A
(Instruction Set Principles) exercises A.1--A.23. A map of every exercise
with its page and relevance to Lecture~2, worked solutions for the ones our
lectures cover, and hints for the rest.}

% ============================================================================
\section{How to Use This File}

\begin{itemize}[leftmargin=*]
  \item The exercise text is \textbf{not} copied here -- open the book
  (\texttt{references/Computer Architecture - A Quantitative Approach.pdf})
  at the page given. ``Book p.'' is the printed page (Appendix~A pages are
  numbered A-1, A-2, \dots); ``PDF p.'' is the page in the PDF viewer.
  \item The book uses 64-bit MIPS (\texttt{LD}, \texttt{SD}, \texttt{DADD},
  \dots); the ideas are the same as the 32-bit MIPS of our lecture.
  \item A.1--A.4 are the CPI-from-instruction-mix exercises (Lecture~1
  skills applied to MIPS). They need Figures~A.27 and A.28 (book pp.~A-41 and
  A-42, PDF pp.~564--565).
\end{itemize}

% ============================================================================
\section{Exercise Map}

\begin{center}
\small
\renewcommand{\arraystretch}{1.15}
\begin{tabular}{@{}lll>{\raggedright\arraybackslash}p{5.2cm}l@{}}
\toprule
Ex. & Book p. & PDF p. & Topic & Priority \\
\midrule
A.1 & A-47 & 570 & Effective CPI, SPECint mix (gap, gcc) & \textbf{High} \\
A.2 & A-47 & 570 & Effective CPI, SPECint mix (gzip, perlbmk) & \textbf{High} \\
A.3 & A-47/48 & 570--571 & Effective CPI with FP, SPECfp mix (lucas, swim) & \textbf{High} \\
A.4 & A-48 & 571 & Effective CPI with FP (applu, art) & \textbf{High} \\
A.5 & A-48 & 571 & Copy propagation on three statements & Medium \\
A.6 & A-48 & 571 & Lab: gcc optimization levels & Skip (lab) \\
A.7 & A-48 & 571 & MIPS and x86 code for an array loop; counts & \textbf{High} \\
A.8 & A-48/49 & 571--572 & Expanding-opcode encoding in 12 bits & \textbf{High} \\
A.9 & A-49 & 572 & Code size for the four ISA styles & \textbf{High} \\
A.10 & A-49 & 572 & Discussion: are more registers better? & \textbf{High} \\
A.11 & A-50 & 573 & Struct size and alignment, 32/64-bit & Medium \\
A.12--A.17 & A-50/51 & 573--574 & Labs: measuring mixes, AVX, gcc ports, energy & Skip (lab) \\
A.18 & A-51/52 & 574--575 & Memory traffic of the four ISA styles & \textbf{High} \\
A.19 & A-52 & 575 & ISA styles with register-indirect addressing (loop) & Medium \\
A.20 & A-52/53 & 575--576 & Choosing offset lengths from usage data & Medium \\
A.21 & A-53/54 & 576--577 & Lab: displacement histograms & Skip (lab) \\
A.22 & A-54 & 577 & Big/little endian, ASCII, misaligned reads & \textbf{High} \\
A.23 & A-54 & 577 & Discussion: ISA for desktop/server/cloud/embedded & Medium \\
\bottomrule
\end{tabular}
\end{center}

% ============================================================================
\section{Worked Solutions: High Priority}

\subsection*{Method for A.1--A.4}
\addcontentsline{toc}{subsection}{Method for A.1--A.4}
Group the figure's rows into the CPI classes given in the exercise, average
the two programs, then $\CPI = \sum F_i\cdot\CPI_i$. Conditional branches
use $0.6\times2.0 + 0.4\times1.5 = 1.8$. Assumptions used below (state them
in an exam): \texttt{jump}, \texttt{call} and \texttt{return} count as
jumps; every other integer row (add, sub, mul, compare, load imm, shifts,
logic, cond move) is an ALU instruction; \texttt{sub FP} counts as FP add;
FP move/compare/cond-move/other count as ``other FP''.

\ex{A.1}{A-47}{570}
\emph{Asks:} the effective CPI using the average of gap and gcc.
\begin{solution}
Loads$+$stores: gap $36.8\%$, gcc $38.3\%$ $\to 37.55\%$. Cond.\ branches
$10.7\%$. Jumps (jump$+$call$+$return): $4.0\%$, $1.9\%$ $\to 2.95\%$. ALU:
the rest, $48.8\%$.
\[
  \CPI = 0.488(1.0) + 0.3755(1.4) + 0.107(1.8) + 0.0295(1.2)
  = 0.488 + 0.526 + 0.193 + 0.035 \approx 1.24 .
\]
\end{solution}

\ex{A.2}{A-47}{570}
\emph{Asks:} the same with the average of gzip and perlbmk.
\begin{solution}
Loads$+$stores $35.05\%$, cond.\ branches $10.95\%$, jumps $2.75\%$, ALU
$51.25\%$:
\[
  \CPI = 0.5125 + 0.3505(1.4) + 0.1095(1.8) + 0.0275(1.2) \approx 1.23 .
\]
\end{solution}

\ex{A.3}{A-47/48}{570--571}
\emph{Asks:} the effective CPI including FP classes, average of lucas and
swim (Fig.~A.28).
\begin{solution}
\begin{tabular}[t]{@{}lccc|c@{}}
Class (CPI) & lucas & swim & avg & contribution \\
\hline
Integer ALU (1.0) & 20.9 & 44.8 & 32.85 & 0.329 \\
Int load/store (1.4) & 14.0 & 10.4 & 12.2 & 0.171 \\
Cond.\ branch (1.8) & 0.6 & 1.3 & 0.95 & 0.017 \\
FP load/store (1.5) & 34.4 & 21.8 & 28.1 & 0.422 \\
FP add/sub (4.0) & 15.8 & 13.7 & 14.75 & 0.590 \\
FP multiply (6.0) & 9.4 & 6.9 & 8.15 & 0.489 \\
FP divide (20.0) & 0 & 0.3 & 0.15 & 0.030 \\
Other FP (2.0) & 5.0 & 0.9 & 2.95 & 0.059 \\
\end{tabular}

$\CPI \approx 2.11$. FP arithmetic, although only $\approx23\%$ of the
instructions, contributes over half of the cycles.
\end{solution}

\ex{A.4}{A-48}{571}
\emph{Asks:} the same with the average of applu and art.
\begin{solution}
\begin{tabular}[t]{@{}lccc|c@{}}
Class (CPI) & applu & art & avg & contribution \\
\hline
Integer ALU (1.0) & 50.4 & 42.1 & 46.25 & 0.463 \\
Int load/store (1.4) & 16.7 & 18.1 & 17.4 & 0.244 \\
Cond.\ branch (1.8) & 2.5 & 11.5 & 7.0 & 0.126 \\
FP load/store (1.5) & 15.6 & 16.5 & 16.05 & 0.241 \\
FP add/sub (4.0) & 5.2 & 4.5 & 4.85 & 0.194 \\
FP multiply (6.0) & 8.6 & 4.1 & 6.35 & 0.381 \\
FP divide (20.0) & 0.3 & 0.6 & 0.45 & 0.090 \\
Other FP (2.0) & 0.7 & 2.4 & 1.55 & 0.031 \\
\end{tabular}

$\CPI \approx 1.77$.
\end{solution}

\ex{A.7}{A-48}{571}
\emph{Asks:} code for \texttt{A[i] = B[i] + C} ($i = 0..100$, 64-bit data,
everything kept in memory between uses) on MIPS and x86: dynamic
instruction count, memory-data references and code size.
\begin{solution}
(a) One possible MIPS64 version (A, B, C, i at 1000, 3000, 5000, 7000):
\begin{lstlisting}
        DADD  R1, R0, R0       ; i = 0
        SD    R1, 7000(R0)
Loop:   LD    R1, 7000(R0)     ; load i
        DSLL  R2, R1, #3       ; i * 8
        LD    R3, 3000(R2)     ; B[i]
        LD    R4, 5000(R0)     ; C
        DADD  R3, R3, R4
        SD    R3, 1000(R2)     ; A[i]
        DADDI R1, R1, #1       ; i++
        SD    R1, 7000(R0)     ; store i
        SLTI  R5, R1, #101     ; i <= 100 ?
        BNEZ  R5, Loop
\end{lstlisting}
Dynamic count $= 2 + 10\times101 = 1012$. Data references $= 1 + 5\times101
= 506$ (3 loads $+$ 2 stores per iteration). Code size $= 12\times4 = 48$
bytes.

(b) x86 can use memory operands and scaled indexing, e.g.\ per iteration:
\texttt{mov rax,[i]; mov rbx,[B+rax*8]; add rbx,[C]; mov [A+rax*8],rbx;
add qword [i],1; cmp qword [i],100; jle Loop}. That is $1 + 7\times101 =
708$ instructions and $1 + 7\times101 = 708$ data references (the
\texttt{add} to memory reads and writes). Code size depends on the encodings
(variable length, roughly 4--10 bytes each): fewer instructions than MIPS,
but more bytes per instruction and more memory traffic.
\end{solution}

\ex{A.8}{A-48/49}{571--572}
\emph{Asks:} with 12-bit instructions and 5-bit register fields, can we
encode given numbers of two-, one- and zero-address instructions?
\begin{solution}
Use expanding opcodes. Two-address: $12-10 = 2$ opcode bits, 4 patterns; 3
are used, 1 is left as an escape. One-address: the escape plus 5 more opcode
bits gives $2^5 = 32$ patterns. Each one-address pattern left unused
escapes to another 5 bits, i.e.\ 32 zero-address instructions.
\begin{enumerate}[label=(\alph*),leftmargin=*]
  \item 3 / 30 / 45: 2 spare one-address patterns $\to 2\times32 = 64 \ge 45$.
  \textbf{Possible.}
  \item 3 / 31 / 35: 1 spare pattern $\to 32 < 35$. \textbf{Not possible.}
  \item 3 two-address and 24 zero-address need one escape pattern
  ($24\le32$), so at most $32-1 = 31$ one-address instructions.
\end{enumerate}
\end{solution}

\ex{A.9}{A-49}{572}
\emph{Asks:} (a) names per instruction and code size for $C = A + B$ on each
ISA style of Fig.~A.2 (8-bit opcode, 64-bit memory address, 6-bit register
field); (b) a longer sequence, marking destroyed operands and overhead.
\begin{solution}
(a)

\noindent\begin{tabular}{@{}lp{6.2cm}p{4.2cm}@{}}
Style & Code (names per instruction) & Size \\
\hline
Stack & Push A (1), Push B (1), Add (0), Pop C (1) & $3\cdot72+8 = 224$ bits $=28$ B \\
Accumulator & Load A, Add B, Store C (1 each) & $3\cdot72 = 216$ bits $=27$ B \\
Reg-memory & Load R1,A; Add R3,R1,B; Store R3,C & $78+84+78 = 240$ bits $=30$ B \\
Load-store & Load R1,A; Load R2,B; Add R3,R1,R2; Store R3,C & $78+78+26+78 = 260$ bits $=32.5$ B \\
\end{tabular}

(b) Write each sequence the same way. Stack and accumulator machines
destroy their operands (the top of stack / accumulator), so reusing $A$,
$C$ or $D$ needs extra loads or pushes -- those are the overhead
instructions and overhead data bytes. Register machines keep $A$, $C$, $D$
in registers and need no reloads, at the cost of longer instructions.
\end{solution}

\ex{A.10}{A-49}{572}
\emph{Asks:} trade-offs of adding more registers to MIPS.
\begin{solution}
For: fewer loads/stores (less spilling), more values kept in registers
across loops and calls, more room for compiler optimizations (register
allocation, unrolling, scheduling).
Against: more bits per register field (bigger instructions or fewer opcode
and immediate bits), a larger and slower register file (more ports, longer
access time, possibly a longer clock or extra pipeline stage), more state to
save on context switches and calls, more power and area, and diminishing
returns once the compiler cannot use them all. Changing the count also
breaks binary compatibility.
\end{solution}

\ex{A.18}{A-51/52}{574--575}
\emph{Asks:} code for \texttt{A=B+C; B=A+C; D=A-B} on accumulator,
memory-memory, stack and load-store machines; mark reloads and operand
passing; count instruction and data bytes for 16-bit and then 64-bit
addresses/data.
\begin{solution}
Approach: 1-byte opcodes, direct addresses (2 bytes, then 8 bytes),
register fields 4 bits (16 registers). Count, for each style, instruction
bytes fetched and data bytes moved. Typical outcome: accumulator and stack
code reload values ($A$, $C$) because their single storage place is
overwritten; memory-memory has short programs but every operand is a
memory access; load-store loads $B$, $C$ once, keeps everything in
registers and stores only $A$, $B$, $D$. With 16-bit addresses the
memory-oriented styles are competitive. With 64-bit addresses every address
in an instruction costs 8 bytes, so load-store (which names memory least
often) has the lowest total traffic. See our ISA review's practice problem
on $D=(A+B)\times C$ for the same method.
\end{solution}

\ex{A.22}{A-54}{577}
\emph{Asks:} store \texttt{434F 4D50 5554 4552} in an aligned double word,
big and little endian, read it as ASCII, and list misaligned 2-byte
(big-endian) and 4-byte (little-endian) reads.
\begin{solution}
The bytes \texttt{43 4F 4D 50 55 54 45 52} are the ASCII codes of
\texttt{C O M P U T E R}.
\begin{enumerate}[label=(\alph*),leftmargin=*]
  \item Big endian (address 0 to 7): \texttt{43 4F 4D 50 55 54 45 52}
  $\to$ ``COMPUTER''.
  \item Little endian: \texttt{52 45 54 55 50 4D 4F 43} $\to$ ``RETUPMOC''.
  \item Misaligned halfwords start at odd addresses 1, 3, 5:
  \texttt{0x4F4D}, \texttt{0x5055}, \texttt{0x5445}.
  \item Misaligned words that fit inside the double word start at addresses
  1, 2, 3. In little endian the value is
  $\text{byte}[a{+}3]\,\text{byte}[a{+}2]\,\text{byte}[a{+}1]\,\text{byte}[a]$:
  \texttt{0x50555445}, \texttt{0x4D505554}, \texttt{0x4F4D5055}.
\end{enumerate}
\end{solution}

% ============================================================================
\section{Worked Solutions: Medium Priority}

\ex{A.5}{A-48}{571}
\emph{Asks:} apply copy propagation to \texttt{A=B+C; B=A+C; D=A-B} until no
operand is a computed value; where does work go up or down?
\begin{solution}
Using the original $B$ and $C$: $A = B + C$; $B' = (B + C) + C = B + 2C$;
$D = A - B' = (B + C) - (B + 2C) = -C$. Statement 2 got \emph{more}
work (two additions instead of one); statement 3 got much \emph{less} (just
a negation). So the same transformation can help or hurt: an optimizing
compiler must evaluate costs case by case, which is why optimization is
hard.
\end{solution}

\ex{A.11}{A-50}{573}
\emph{Asks:} size of a given C struct on 32- and 64-bit machines, and the
minimum size after reordering its members.
\begin{solution}
Assume each type is aligned to its own size (double 8, pointers 4 or 8).
32-bit: the members land at offsets 0, 1, 4, 8, 16, 20, 24, 32, 36, 40;
the data ends at 44 and the struct is padded to a multiple of 8: \textbf{48
bytes}. Sorted largest first (doubles, then 4-byte items, short, chars):
\textbf{40 bytes}.
64-bit (8-byte pointers): offsets 0, 1, 4, 8, 16, 20, 24, 32, 40, 48; ends
at 52, padded to \textbf{56 bytes}. Reordered: \textbf{48 bytes}.
\end{solution}

\ex{A.20}{A-52/53}{575--576}
\emph{Asks:} average instruction length if offsets may be 0, 8, 16 or 24
bits, versus a fixed 24-bit instruction (with extra instructions for long
offsets) or always-24-bit offsets.
\begin{solution}
Method: an $n$-bit offset including sign holds $n-1$ magnitude bits, so read
the cumulative percentages of Fig.~A.31 at 0, 7, 15 and 23 bits (data:
30.4\%, 66.9\%, 100\%; branches: 0.1\%, 85.2\%, 98.5\%, 100\%). Weight each
offset length by the load/store and branch frequencies of Fig.~A.27 (ALU
and other instructions are 16 bits), and compare the bytes fetched. The
point of the exercise: variable-length offsets save bytes on average, but a
fixed format is simpler to decode, which is the trade-off behind MIPS's
fixed 16-bit immediates.
\end{solution}

\ex{A.19, A.23}{A-52, A-54}{575, 577}
\begin{solution}
A.19: same method as A.7 and A.18, writing the loop for each ISA style;
register-indirect addressing lets the stack/accumulator machines walk the
arrays with a pointer instead of absolute addresses. A.23: desktop/server
favour performance and compatibility (rich ISAs, 64-bit, SIMD);
cloud/warehouse favours throughput per watt and virtualization support;
embedded favours code density, low power and cost (compressed encodings
such as Thumb/MIPS16, small register files).
\end{solution}

\end{document}
